% This is the ADASS_template.tex LaTeX file, 24 Oct 2025.
% It is based on the ASP general author template file, but modified to reflect the specific
% requirements of the ADASS proceedings.

% This is the aspauthor.tex LaTeX file
% Copyright 2021, Astronomical Society of the Pacific Conference Series
% Revision:  29 November 2021

% To compile, at the command line positioned at this folder, type:
% pdflatex ADASS_template
% pdflatex ADASS_template
% This will create a file called ADASS_template.pdf

\documentclass[11pt,twoside]{article}

% Do NOT use ANY packages other than asp2023. 
\usepackage{asp2023}

\aspSuppressVolSlug
\resetcounters

% References must all use BibTeX entries in a .bibfile.
% References must be cited in the text using \citet{} or \citep{}.
% Do not use \cite{}.
% See ManuscriptInstructions.pdf for more details
\bibliographystyle{asp2021}

% The ``markboth'' line sets up the running heads for the paper.
% 1 author: "Surname"
% 2 authors: "Surname1 and Surname2"
% 3 authors: "Surname1, Surname2, and Surname3"
% >3 authors: "Surname1 et al."
% Replace ``Short Title'' with the actual paper title, shortened if necessary.
% Use mixed case type for the shortened title
% Ensure shortened title does not cause an overfull hbox LaTeX error
% See ASPmanual2010.pdf 2.1.4  and ManuscriptInstructions.pdf for more details
\markboth{Ruta Kale}{Towards reliable batch processing}

\begin{document}

\title{Towards Reliable Batch Data Processing for the SKA Science Data Processor[DRAFT]}

% Note the position of the comma between the author name and the 
% affiliation number.
% Authors surnames should come after first names or initials, eg John Smith, or J. Smith.
% Author names should be separated by commas.
% The final author should be preceded by "and".
% Affiliations should not be repeated across multiple \affil commands. If several
% authors share an affiliation this should be in a single \affil which can then
% be referenced for several author names. If only one affiliation, no footnotes are needed.
% See ManuscriptInstructions.pdf and ASP's manual2010.pdf 3.1.4 for more details
\author{Ruta~Kale,$^{1,2}$}
\affil{$^1$SKA Observatory, Jodrell Bank, Lower Withington, Macclesfield, Cheshire, SK11 9FT, United Kingdom; \email{ruta.kale@skao.int}}
\affil{$^2$National Centre for Radio Astrophysics, Tata Institute of Fundamental Research, Savitribai Phule Pune University Campus, Ganeshkhind, Pune 411007 Maharashtra, India}

% This section is for ADS Processing.  There must be one line per author. paperauthor has 9 arguments.
\paperauthor{Ruta~Kale}{ruta.kale@skao.int}{https://orcid.org/0000-0003-1449-3718}{SKA Observatory}{Author1 Department}{City}{Cheshire}{SK119FT}{United Kingdom}
%\paperauthor{Sample~Author2}{Author2Email@email.edu}{ORCID_Or_Blank}{Author2 Institution}{Author2 Department}{City}{State/Province}{Postal Code}{Country}
%\paperauthor{Sample~Author3}{Author3Email@email.edu}{ORCID_Or_Blank}{Author3 Institution}{Author3 Department}{City}{State/Province}{Postal Code}{Country}

% There should be one \aindex line (commented out) for each author. These are used to
% build up the author index for the Proceedings. The surname must come first, followed by
% initials. Note the use of ~ before each initial to control spacing.
% The \author entries (see above) have surname last. These \aindex entries have
% surname first.
% The Aindex.py command willl create them for you after you have constructed the \author
% The first entry should be the first author, for bold-facing the author index page-reference

%\aindex{FistAuthor1,~S.~A.}
%\aindex{Author2,~S.~B.}
%\aindex{Author3,~S.}


\begin{abstract}
The SKA Science Data Processor (SDP) must turn large visibility datasets
into scientifically validated images within finite computing and storage
resources. We describe the operational context, processing scale and
emerging batch-workflow design. Illustrative data volumes and image
dimensions show why memory use and data movement matter alongside compute.
Modular scientific stages, explicit interfaces and Measurement Set version
4 (MSv4) support evolving calibration and imaging methods. Reproducible
configuration, provenance and recoverable execution connect this design
to delivery from commissioning through science operations. We identify
validation priorities for scientific correctness, restart behaviour and
throughput; these are design objectives, rather than measured pipeline
performance or demonstrated production readiness.
\end{abstract}

% These lines show examples of subject index entries. At this stage these have to commented
% out, and need to be on separate lines. Eventually, they will be automatically uncommented
% and used to generate entries in the Subject Index at the end of the Proceedings volume.
% Don't leave these in! - replace them with ones relevant to your paper.
%\ssindex{FOOBAR!conference!ADASS 2019}
%\ssindex{FOOBAR!organisations!ASP}

% These lines show examples of ASCL index entries. At this stage these have to commented
% out, and need to be on separate lines. Eventually, they will be automatically uncommented
% and used to generate entries in the ASCL Index at the end of the Proceedings volume.
% The ascl.py command will scan your paper on possible code names.
% Don't leave these in! - replace them with ones relevant to your paper.
%\ooindex{FOOBAR, ascl:1101.010}

\section{Introduction}
The Square Kilometre Array Observatory (SKAO) is constructing SKA-Low in
Western Australia (50--350 MHz) and SKA-Mid in the South African Karoo
(0.35--15.4 GHz). The SKA-Low is an aperture array\footnote{256 individual antenna elements together form a station, and the telescope is an array of stations.} and the SKA-Mid is an array of dishes.
The Array Assemblies (AAs) provide successive configurations of the telescopes as the construction progresses:
Low-AA2 has 64 Low stations and Mid-AA2 has 64 dishes; the construction scope comprises
257 Low stations at AA$^*$ (Mod) and 144 Mid dishes at AA$^*$. The full
AA4 design has 512 stations and 197 dishes, but is not a committed
construction stage. Mid AA$^*$ and AA4 include 64 MeerKAT
dishes\footnote{\raggedright\url{https://www.skao.int/en/science-users/118/ska-telescope-specifications}}.

Interferometers record complex correlations, or visibilities. Calibration
estimates instrumental and propagation effects; imaging reconstructs sky
brightness from incomplete spatial-frequency sampling. The broad scope of Science Data Processor starts from getting a request for observations from the user, processing it to make the observations, perform real time and batch processing and deliver the data products to the user. The aspect on which is contribution focuses is batch processing - the part where the SDP performs calibration and imaging
and manages the resulting data products. Reliable batch processing must
 preserve scientific validity while capabilities and algorithms
evolve. Following \citet{Rees2019}, we treat software quality as both a
scientific and an operational requirement.

\section{Operational Context and Scale}
Users specify products and processing requirements during proposal submission and if the proposal is successful, the observations are carried out, the data are processed and data products are delivered. All of these steps are orchestrated through the SDP. The SDP creates observation-level products for
quality assessment (QA) and delivery to the SKA Regional Centre Network
(SRCNet), where users perform further analysis 
%Routine raw and intermediate data are planned to be deleted after all
%planned processing that needs them has completed
\footnote{\raggedright
SKAO Science Data Products: A Summary, revision 02 (2025):
\url{https://www.skao.int/en/science-users/649/skao-science-data-products}}.
Provenance and QA must therefore support confidence in products before
their inputs are discarded. Our focus is imaging, rather than all SDP modes.

For $N$ stations or dishes, the number of cross-correlation baselines is
$N_{\rm bl}=N(N-1)/2$, excluding autocorrelations. The visibility payload is
$V_{\rm vis}=N_{\rm bl}N_tN_{\rm chan}N_{\rm pol}b$, where the remaining
factors count time samples, channels, correlation products and bytes per
complex value. For one hour, one-second sampling, 4,096 channels, four
products and eight bytes per complex value, AA2 has 2,016 baselines and
0.95 TB for either telescope. Low AA$^*$ (Mod) has 32,896 baselines and
15.52 TB; Mid AA$^*$ has 10,296 baselines and 4.86 TB. These represent
16.32- and 5.11-fold increases over AA2. Full-design AA4 gives 130,816
baselines and 61.73 TB for Low, and 19,306 baselines and 9.11 TB for Mid.
These decimal-terabyte payloads exclude flags, weights, metadata and
storage overhead and thus are minimal estimates; actual volumes depend on sampling, averaging and duration.

Simplistically adopting the field of view as $\theta_{\rm FoV}=\lambda/D$ and resolution as
$\theta_{\rm res}=\lambda/B_{\max}$, and the cellsize as 
$\theta_{\rm cell}=\theta_{\rm res}/3$, an image covering the field at the
band minimum with cells set at the band maximum has width
\begin{equation}
 N_{\rm pix}=\frac{3B_{\max}}{D}\frac{\nu_{\max}}{\nu_{\min}}.
 \label{eq:imagewidth}
\end{equation}
The D denotes station or dish diameter and $B_{max}$ denotes the longest baseline. The wavelength $\lambda$ for the field of view calculation is corresponding to the lowest frequency of the observing band and for the resolution calculation is considered to be the one corresponding to the highest frequency of the band to underline the effect of the wide bandwidths. These are taken for illustrative purpose and will depend in practice on the observing parameters and the science goal.
For Low, $D=39$ m and adopted maximum baselines of 39 km (AA2) and
73.4 km (AA$^*$ (Mod)/AA4) give widths of 21,000 and approximately
39,500 pixels across 50--350 MHz.  Resolution goes from 31.71 to
4.53 arcsec across the band for AA2, and from 16.85 to 2.41 arcsec for
AA$^*$ (Mod). For Mid, $D=15$ m and the adopted AA2/AA$^*$ baseline is
108 km including SKA008, or 36 km without it. Band 1 (0.35--1.05 GHz)
then gives widths of 64,800 or 21,600 pixels, respectively, and upper-edge
resolutions of 0.55 or 1.64 arcsec. Thus the adopted Mid resolution and
image width do not change between AA2 and AA$^*$. Mid AA4 uses
159.6 km, unaffected by
SKA008, and gives 95,760 pixels: about 9.2 billion pixels per square image
plane. Band 2 (0.95--1.76 GHz) gives approximately 40,000, 13,300 and
59,100 pixels for these three Mid cases. 
%MeerKAT's 13.5 m dishes have
%a different primary beam. These are sizing illustrations; science goals,
%source distribution, bandwidth, weighting, padding and faceting determine
%the operational image layout.

% AUTHOR CHECK: confirm the approved hardware specification and publication permission.
The SDP hardware for SKA-Low is in the phase of testing a scalable unit.
The supplied Low AA2 scalable-unit has 124 dual-socket nodes,
two Intel Xeon 6 6972P processors per node and 23,808 CPU cores.
Its 768 GB per node totals approximately 95 TB, with about 1.75 PFLOP/s
peak FP64 and 10 PB raw DDN Lustre storage from 415 24-TB disks.
Four network fabrics and a shared core support expansion are planned. These are
platform specifications, not pipeline benchmarks: usable storage is lower,
and throughput depends on memory access, I/O and algorithm efficiency. The batch processing needs to work within the constraints of these specifications.

\section{Workflow Interfaces and Recovery}
Batch stages include preprocessing, instrumental calibration, iterative
self-calibration, and continuum or spectral-line imaging. Their ordering
follows the scientific recipe for the planned data products.
Figure~\ref{fig:architecture} separates scientific intent, workflow state
and resource execution. Recipes specify operations and dependencies.
Workflow control tracks readiness and requests resources through execution adapters.
Resolved configurations identify data selection, algorithm parameters,
execution settings and defaults. Interfaces must specify scientific inputs,
solution validity, partitioning assumptions, output schemas and completion
criteria. Dataset locations alone do not establish whether data are observed
or corrected, which flags and weights apply, or whether a product is complete. The batch processing needs to be scalable and current algorithms have limitations. Thus the pipelines need to have the flexibility to adapt new algorithms as they are developed. 

\begin{figure}[ht]
\centering
\setlength{\unitlength}{1mm}
\footnotesize
\begin{picture}(130,43)
\put(1,30){\framebox(41,11){\shortstack{Recipe and parameters\\Scientific intent}}}
\put(49,30){\framebox(33,11){\shortstack{Workflow control\\Dependencies and state}}}
\put(89,30){\framebox(40,11){\shortstack{Execution adapters\\Compute and storage}}}
\put(42,35){\vector(1,0){7}}
\put(82,35){\vector(1,0){7}}
\put(65,30){\vector(0,-1){7}}
\put(49,12){\framebox(33,11){\shortstack{Scientific stages\\Calibration / imaging}}}
\put(89,12){\framebox(40,11){\shortstack{Data interfaces\\MSv4 / stage products}}}
\put(82,17){\vector(1,0){7}}
\put(1,0){\framebox(128,7){Proposed shared records: provenance, status, logs and quality assessment}}
\end{picture}
\caption{Conceptual batch framework. Arrows indicate coordination and
interfaces, not visibility-data transport or deployed-service status.}
\label{fig:architecture}
\end{figure}

There are several limitations of the "MSv2" format of the visibilities that is currently in use. Thus a new scalable format, MSv4, is being developed through a global collaboration.
MSv4 represents visibilities as labelled arrays with time, baseline,
frequency and polarisation coordinates; related datasets form a Processing
Set. XRADIO supplies implementation, access and conversion
tools\footnote{\raggedright
\url{https://xradio.readthedocs.io/en/latest/measurement_set/overview.html}},
through collaboration among CASA partners NRAO, ESO, NAOJ and JIVE,
SKAO and SARAO. XRADIO supports selective, lazy loading and distributed
processing; Zarr provides chunked, compressed arrays and object-storage
access. 
%Chunk sizes must match algorithm access patterns and working sets.
%MSv2 conversion requires checks on coordinates, flags, weights and metadata
%because the models are not backward compatible. These capabilities do not
%establish production readiness: batch-preprocessing version 2.7.3 describes
%its XRADIO path as experimental and outside the end-to-end
%pipeline\footnote{\raggedright
%\url{https://developer.skao.int/projects/ska-sdp-batch-preprocess/en/2.7.3/xradio_pipeline.html}}.

The batch pipelines in addition need to be resilient.
Validated stage outputs provide recovery boundaries, avoiding a complete
restart after failure. Retries must identify inputs and configuration,
protect completed products and reject partial outputs. Checkpoint placement
balances recomputation against additional writes and storage. Provenance
must link input identifiers, software versions, resolved configuration,
calibration products and model history to workflow state, logs and QA.
Reproducibility requires scientific tolerances, rather than necessarily
bitwise equality across hardware or parallel reduction orders. These are
design requirements, not evidence that every recovery path is implemented.

\section{Delivery and Validation}
The batch pipelines need to satisfy requirements ranging from science commissioning of the telescopes, the science verification phase and through the science operations.
Commissioning needs rapid configuration changes, diagnostics and repeatable
reprocessing. Production workflows can be useful for the global sky model survey. Low AA2 science verification adds the requirement of production workflows for available observing modes and throughput matched to the observing plan. Science operations broadens the supported modes which the batch pipelines need to be ready for. The framework of batch pipeline  must also accommodate
algorithm evolution and both telescopes throughout this progression.

Validation must test four aspects together: scientific correctness through
reference comparisons of flux recovery, noise and spectral fidelity; safe
recovery through failure injection and restart tests; platform fit through
wall time, peak memory, I/O and storage measurements; and scaling through
representative observations and multi-node runs under realistic concurrency.
More workers may expose contention as well as parallelism. A successful run
is necessary; repeatable science and sustained throughput establish readiness.

\section{Conclusions}
The SKAO SDP is set to deliver reliable batch processing through the phases of science commissioning, science verification and science operations. The reproducible configuration, provenance and recoverable execution connect the design to delivery from commissioning through science operations.

%Visibility volumes and image dimensions make memory use and data movement
%central to SDP batch processing. Modular stages, explicit MSv4 interfaces
%and separate execution control allow algorithms to evolve within a stable
%framework. Reproducible configuration, provenance and recoverable workflows
%must be tested on representative workloads. The goal is trusted,
%scientifically validated products from commissioning through science
%operations as the telescopes and processing methods evolve.

%\bibliography{references}

%%%%%%%%%%%


\acknowledgements We recognise and acknowledge the Indigenous peoples and cultures that have traditionally lived on the lands on which our facilities are located. 


\bibliography{C409}  % For BibTex

% if we have space left, we might add a conference photograph here. Leave commented for now.
% \bookpartphoto[width=1.0\textwidth]{foobar.eps}{FooBar Photo (Photo: Any Photographer)}

\end{document}
